\clearpage\begingroup
\HistoricalHead{185}{R}{WECHSELWIRKUNGEN ZWISCHEN GALVANISCHEN STRÖMEN UND MAGNETISMUS.}
\begin{center}7.\end{center}
\noindent\hspace*{\parindent}Dieselben Formeln werden auch noch hinlängliche Genauigkeit geben, wenn wir Statt der getrennten Ströme in $n$ Rechtecken einen zusammenhängenden annehmen, der sich gleichsam Schraubenförmig so um das Prisma schlingt, dass die erste Umwindung zwischen dem ersten und zweiten
\HistoricalFoot{XI$_1$.}{24}
\clearpage\SourceLeaf{191}\HistoricalHead{186}{L}{PHYSIK.\quad NACHLASS.}
\noindent Rechteck, die zweite zwischen dem zweiten und dritten Rechteck liegt u. s. w. Die Anzahl dieser Umwindungen wird nun $n-1$ und die Integration schlechthin zwischen den äussersten Werthen von $k$ auszuführen sein, deren Differenz $\lambda(n-1)$ sein wird. Eine grössere Genauigkeit in diese Berechnung zu legen würde ohne Nutzen, obwohl ohne weitere Schwierigkeit als grössere Weitläuftigkeit des Calcüls sein.

Auf diese Weise lässt sich also die Kraft, mit welcher ein einen Multiplicator durchlaufender galvanischer Strom auf das magnetische Fluidum in jedem Punkte des Raumes wirkt, berechnen. Wir wollen diese Berechnung an einigen Beispielen erläutern.
\clearpage\endgroup
